Equipes de software combinam \emph{sistemas puxados} ao estilo Kanban com cadências Scrum, porém modelos quantitativos explícitos costumam omitir relações interpessoais (sinergia) que afetam colaboração e repasses. Este artigo descreve um modelo operacional de simulação \textbf{em tempo discreto} (granularidade diárias) com colunas de fluxo, capacidade estocástica por membro, políticas de WIP e planejamento, e multiplicadores de sinergia para trabalho colaborativo e para repasses entre etapas, incluindo retrabalho. O motor é reprodutível via semente de PRNG. Relatamos um conjunto mínimo de cenários comparativos (A/B) com a mesma semente e \emph{backlog}, variando apenas sinergia entre pessoas, juntamente com métricas de fluxo (contagens por coluna ao longo do tempo, tempo de ciclo e \emph{throughput} por sprint). A contribuição é metodológica: um arcabouço explícito para análise de sensibilidade e ensino em pesquisa operacional, com fronteira clara entre hipóteses operacionais e validação empírica.
